% ==============================================================================
\section{Experimental Setup}
\label{sec:setup}
% ==============================================================================

To evaluate the performance and scalability of our parallel I/O strategies, we conducted experiments across two different High-Performance Computing (HPC) clusters: a regional cluster used for initial setup and prototyping, and a national supercomputer used for larger-scale runs.

\subsection{Hardware and Storage Infrastructure}

The first phase of our experiments took place on the Plafrim cluster at Inria Bordeaux. Benchmarks were deployed on the Bora compute nodes, which are equipped with dual-socket Intel Xeon Gold 6240 processors (36 physical cores per node at 2.60~GHz) and 192~GB of RAM per node. 
Inter-node communications relied on an Intel Omni-Path network delivering 100~Gb/s bandwidth. Storage on Plafrim is managed by a BeeGFS parallel file system offering an aggregate bandwidth capped at 10~Gbit/s and configured with up to 8 active Object Storage Targets (OSTs). This platform provided a controlled environment to validate our 
setup and measure baseline performance.
To evaluate our strategies at a larger scale, we extended our benchmark campaign to the Adastra supercomputer at CINES. Specifically, tests were run on the GENOA CPU partition, which features dual-socket AMD EPYC 9654 processors with 192 physical cores per node and high memory throughput. Nodes are connected via an HPE Slingshot-11 high-speed network. 
Storage on Adastra is provided by a production-grade Lustre parallel file system, which allows fine-grained data striping across up to 40 active OSTs.

\subsection{Benchmark Workload \& Evaluation Parameters}

The experiments considered several 5D grid configurations representative of different Gysela workloads. The evaluated resolutions and their corresponding single-checkpoint memory footprints are summarized below:

\begin{align*}
    C_1 &: N_r \times N_\theta \times N_\varphi \times N_{v_\parallel} \times N_\mu
        = 128 \times 128 \times 64 \times 64 \times 64 \\
    C_2 &: N_r \times N_\theta \times N_\varphi \times N_{v_\parallel} \times N_\mu
        = 256 \times 512 \times 128 \times 128 \times 64 \\
    C_3 &: N_r \times N_\theta \times N_\varphi \times N_{v_\parallel} \times N_\mu
        = 512 \times 1024 \times 64 \times 128 \times 64
\end{align*}

These configurations result in single-checkpoint memory footprints ranging from approximately 64 GB to \dots GB. Each configuration was evaluated using the different process topologies considered in the experiments, allowing their impact on I/O performance to be assessed independently of the problem size.

Benchmarks were submitted in exclusive node allocation mode with 8 tasks per node, across process scales ranging from 16 MPI ranks (2 nodes) up to  MPI ranks 2048 (256 nodes). To account for background network and storage noise, each experimental point represents the average of 5 independent executions.

